\section{Forced Plate Equation Python Implementation}

\begin{lstlisting}{language=Python}
import numpy as np
import matplotlib.pyplot as plt
from matplotlib.figure import Figure
from scipy import sparse
from scipy.sparse.linalg import splu
from typing import Callable

def _apply_free_bc(W: np.ndarray, nu: float) -> np.ndarray:

    W[0, 1:-1] = (2.0*W[1, 1:-1] - W[2, 1:-1]
                  - nu*(W[1, 2:] - 2.0*W[1, 1:-1] + W[1, :-2]))
    # Right ghost (row -1), edge is row -2, next interior is row -3
    W[-1, 1:-1] = (2.0*W[-2, 1:-1] - W[-3, 1:-1]
                   - nu*(W[-2, 2:] - 2.0*W[-2, 1:-1] + W[-2, :-2]))
    # Bottom ghost (col 0)
    W[1:-1, 0] = (2.0*W[1:-1, 1] - W[1:-1, 2]
                  - nu*(W[2:, 1] - 2.0*W[1:-1, 1] + W[:-2, 1]))
    # Top ghost (col -1)
    W[1:-1, -1] = (2.0*W[1:-1, -2] - W[1:-1, -3]
                   - nu*(W[2:, -2] - 2.0*W[1:-1, -2] + W[:-2, -2]))
    # Corner condition
    W[0,  0]  = W[2,  0]  + W[0,  2]  - W[2,  2]
    W[-1, 0]  = W[-3, 0]  + W[-1, 2]  - W[-3, 2]
    W[0,  -1] = W[2,  -1] + W[0,  -3] - W[2,  -3]
    W[-1, -1] = W[-3, -1] + W[-1, -3] - W[-3, -3]

    return W

def _laplacian_padded(W: np.ndarray, h: float) -> np.ndarray:
    return (
        W[:-2, 1:-1] + W[2:, 1:-1]
        + W[1:-1, :-2] + W[1:-1, 2:]
        - 4.0 * W[1:-1, 1:-1]
    ) / h ** 2

def _biharmonic(w_phys: np.ndarray, h: float, nu: float) -> np.ndarray:
    N = w_phys.shape[0] - 1

    W = np.zeros((N + 3, N + 3))
    W[1:-1, 1:-1] = w_phys
    W = _apply_free_bc(W, nu)

    lap1 = _laplacian_padded(W, h)

    L = np.zeros((N + 3, N + 3))
    L[1:-1, 1:-1] = lap1
    L = _apply_free_bc(L, nu)

    return _laplacian_padded(L, h)


def _laplacian_phys(w_phys: np.ndarray, h: float, nu: float) -> np.ndarray:
    N = w_phys.shape[0] - 1
    W = np.zeros((N + 3, N + 3))
    W[1:-1, 1:-1] = w_phys
    W = _apply_free_bc(W, nu)
    return _laplacian_padded(W, h)


def _Kh(w: np.ndarray, h: float, nu: float,
        K0: float, T: float, D: float) -> np.ndarray:
    out = K0 * w
    if T != 0.0:
        out = out - T * _laplacian_phys(w, h, nu)
    if D != 0.0:
        out = out + D * _biharmonic(w, h, nu)
    return out


def _Bh(w: np.ndarray, h: float, nu: float,
        K1: float, T1: float) -> np.ndarray:
    out = K1 * w
    if T1 != 0.0:
        out = out - T1 * _laplacian_phys(w, h, nu)
    return out


def _build_system_matrix(
    N: int, h: float, nu: float,
    rho_h: float, K0: float, T: float, D: float,
    K1: float, T1: float,
    beta: float, gamma: float, dt: float,
) -> object:

    n_pts = (N + 1) ** 2

    cols, rows, vals = [], [], []
    e = np.zeros((N + 1, N + 1))

    for col_idx in range(n_pts):
        i, j = divmod(col_idx, N + 1)
        e[i, j] = 1.0

        col_vec = (
            rho_h * e
            + beta * dt ** 2 * _Kh(e, h, nu, K0, T, D)
            + gamma * dt     * _Bh(e, h, nu, K1, T1)
        )

        nz = np.flatnonzero(col_vec)
        for row_idx in nz:
            rows.append(row_idx)
            cols.append(col_idx)
            vals.append(col_vec.flat[row_idx])

        e[i, j] = 0.0

    M = sparse.csc_matrix((vals, (rows, cols)), shape=(n_pts, n_pts))
    return M


def solve_kirchhoff_love(
    force_fn: Callable[[np.ndarray, np.ndarray, float], np.ndarray],
    L: float = 0.24,
    N: int = 80,
    T: float = 1.0,
    rho_h: float = 2.7,
    E: float = 69e9,
    h_plate: float = 0.001,
    nu: float = 0.33,
    K0: float = 0.0,
    T_tension: float = 0.0,
    K1: float = 0.0,
    T1: float = 0.0,
    dt: float | None = None,
    save_path: str | None = None,
):
    D = E * h_plate ** 3 / (12.0 * (1.0 - nu ** 2))
    h = L / N
    if dt is None:
        dt = 1 * h / 2.0

    beta_nb  = 0.25
    gamma_nb = 0.50

    x = np.linspace(0.0, L, N + 1)
    y = np.linspace(0.0, L, N + 1)
    X, Y = np.meshgrid(x, y, indexing="ij")

    n_pts = (N + 1) ** 2

    w = np.zeros((N + 1, N + 1))   # displacement
    v = np.zeros((N + 1, N + 1))   # velocity

    F0 = force_fn(X, Y, 0.0)
    a  = (F0 - _Kh(w, h, nu, K0, T_tension, D)
              - _Bh(v, h, nu, K1, T1)) / rho_h

    print("Building and factorising system matrix ", end="", flush=True)
    M_sys = _build_system_matrix(
        N, h, nu, rho_h, K0, T_tension, D, K1, T1,
        beta_nb, gamma_nb, dt
    )
    lu = splu(M_sys)
    print("done.")

    total_steps = max(1, int(round(T / dt)))
    t_current   = 0.0

    ts = []
    ws = []
    for step in range(total_steps):
        t_new = t_current + dt

        # Stage 1: predictor
        w_p = w + dt * v + 0.5 * dt ** 2 * (1.0 - 2.0 * beta_nb) * a
        v_p = v + dt * (1.0 - gamma_nb) * a

        # Stage 2: solve for a_new
        F_new = force_fn(X, Y, t_new)
        rhs   = F_new - _Kh(w_p, h, nu, K0, T_tension, D) \
                      - _Bh(v_p, h, nu, K1, T1)

        a_new_flat = lu.solve(rhs.ravel())
        a_new = a_new_flat.reshape(N + 1, N + 1)

        # Stage 3: corrector
        w = w_p + beta_nb  * dt ** 2 * a_new
        v = v_p + gamma_nb * dt      * a_new
        a = a_new

        t_current = t_new

        if (step + 1) % max(1, total_steps // 10) == 0:
            pct = 100 * (step + 1) / total_steps
            print(f"  {pct:.0f}% complete (t = {t_current:.4f} s)")
            ts.append(t_current)
            ws.append(w)
    ts.append(t_current)
    ws.append(w)
    return ts,ws
\end{lstlisting}